\documentclass[11pt,a4paper]{article}

\usepackage[T1]{fontenc}
\usepackage[utf8]{inputenc}
\usepackage[ngerman]{babel}
\usepackage{lmodern}
\usepackage{microtype}
\usepackage[a4paper,margin=2.5cm]{geometry}

\usepackage{amsmath,amssymb,amsthm,mathtools}
\usepackage{booktabs}
\usepackage{enumitem}
\usepackage{hyperref}

\hypersetup{
	colorlinks=true,
	linkcolor=blue,
	citecolor=blue,
	urlcolor=blue,
	pdftitle={Ein diskreter Vierlingsoperator auf dem neutralen Sektor der Klein-Vierergruppe},
	pdfauthor={Thomas Hoffbauer}
}

\title{Ein diskreter Vierlingsoperator auf dem neutralen Sektor der Klein-Vierergruppe}
\author{Thomas Hoffbauer}
\date{\today}

\newtheorem{definition}{Definition}[section]
\newtheorem{lemma}[definition]{Lemma}
\newtheorem{proposition}[definition]{Proposition}
\newtheorem{corollary}[definition]{Korollar}
\theoremstyle{remark}
\newtheorem{remark}[definition]{Bemerkung}

\newcommand{\V}{V_4}
\newcommand{\X}{\mathcal X}
\newcommand{\XE}{\mathcal X_E}
\newcommand{\HE}{\mathcal H_E}
\newcommand{\wt}{\mathrm{wt}}

\begin{document}
	\maketitle
	
	\begin{abstract}
		Es wird ein endlicher Operatoransatz auf dem neutralen Vierlingssektor der Klein-Vierergruppe formuliert. Der Zustandsraum besteht aus allen Vierlingen
		\[
		(\sigma_1,\sigma_2,\sigma_3,\sigma_4)\in \V^4
		\]
		mit neutralem Gesamtprodukt
		\[
		\sigma_1\sigma_2\sigma_3\sigma_4=E.
		\]
		Auf dem zugehörigen komplexen Vektorraum wird ein selbstadjungierter Operator aus zyklischem Shift, klassenerhaltenden Paarflips und einem diagonalen Potential definiert. Der Zustandsraum besitzt eine kanonische Zerlegung in Shift-Orbits. Numerisch zeigt der Grundzustand für eine feste Parameterwahl einen Wechsel zwischen zwei Typklassen neutraler Vierlinge, nämlich ABCE- und AABB-Konfigurationen. Der Umschlag wird durch einen Potentialparameter \(\alpha\) gesteuert und liegt numerisch bei ungefähr \(\alpha_c\approx -0.72\).
	\end{abstract}
	
	\section{Grundkonstruktion}
	
	\begin{definition}
		Sei
		\[
		\V=\{E,A,B,C\}
		\]
		die Klein-Vierergruppe mit Multiplikation
		\[
		A^2=B^2=C^2=E,\qquad AB=C,\qquad AC=B,\qquad BC=A.
		\]
	\end{definition}
	
	\begin{definition}
		Der volle Vierlingsraum ist
		\[
		\X=\V^4.
		\]
		Für
		\[
		X=(\sigma_1,\sigma_2,\sigma_3,\sigma_4)\in \X
		\]
		definieren wir die Gesamtsignatur
		\[
		Q(X):=\sigma_1\sigma_2\sigma_3\sigma_4\in \V.
		\]
	\end{definition}
	
	\begin{definition}
		Der neutrale Vierlingssektor ist
		\[
		\XE:=\{X\in \X : Q(X)=E\}.
		\]
		Der zugehörige komplexe Zustandsraum ist
		\[
		\HE:=\mathbb C[\XE].
		\]
	\end{definition}
	
	\begin{lemma}
		Es gilt
		\[
		|\XE|=64.
		\]
	\end{lemma}
	
	\begin{proof}
		Wählt man \(\sigma_1,\sigma_2,\sigma_3\in \V\) beliebig, so ist
		\[
		\sigma_4=\sigma_1\sigma_2\sigma_3
		\]
		eindeutig bestimmt und erfüllt
		\[
		\sigma_1\sigma_2\sigma_3\sigma_4=E,
		\]
		da jedes Element von \(\V\) sein eigenes Inverses ist. Also ist \(|\XE|=4^3=64\).
	\end{proof}
	
	\begin{lemma}
		Jede Permutation der vier Koordinaten erhält \(Q\).
	\end{lemma}
	
	\begin{proof}
		Die Gruppe \(\V\) ist abelsch.
	\end{proof}
	
	\section{Operatoren auf dem neutralen Sektor}
	
	\begin{definition}
		Der Shiftoperator \(T\) auf \(\HE\) ist auf Basiszuständen durch
		\[
		T|\sigma_1,\sigma_2,\sigma_3,\sigma_4\rangle
		=
		|\sigma_2,\sigma_3,\sigma_4,\sigma_1\rangle
		\]
		definiert.
	\end{definition}
	
	\begin{lemma}
		Der Operator \(T\) erhält \(\XE\) und erfüllt \(T^4=I\).
	\end{lemma}
	
	\begin{proof}
		Die Erhaltung von \(\XE\) folgt aus der Permutationsinvarianz von \(Q\). Die Gleichung \(T^4=I\) ist unmittelbar.
	\end{proof}
	
	\begin{definition}
		Für \(g\in\{A,B,C\}\) und \(1\le i<j\le 4\) sei
		\[
		F_{ij}^{(g)}
		|\sigma_1,\dots,\sigma_i,\dots,\sigma_j,\dots,\sigma_4\rangle
		=
		|\sigma_1,\dots,g\sigma_i,\dots,g\sigma_j,\dots,\sigma_4\rangle.
		\]
	\end{definition}
	
	\begin{lemma}
		Jeder Operator \(F_{ij}^{(g)}\) erhält \(\XE\).
	\end{lemma}
	
	\begin{proof}
		Da \(g^2=E\), gilt
		\[
		(g\sigma_i)(g\sigma_j)=g^2\sigma_i\sigma_j=\sigma_i\sigma_j.
		\]
		Somit bleibt das Produkt aller vier Koordinaten unverändert.
	\end{proof}
	
	\begin{definition}
		Sei \(v:\XE\to\mathbb R\) eine Funktion. Der zugehörige diagonale Operator \(V\) auf \(\HE\) sei gegeben durch
		\[
		V|X\rangle=v(X)|X\rangle.
		\]
	\end{definition}
	
	\begin{definition}
		Für Parameter \(t,\lambda\ge 0\) definieren wir
		\[
		H
		=
		-t(T+T^{-1})
		-\lambda\sum_{1\le i<j\le 4}\sum_{g\in\{A,B,C\}}F_{ij}^{(g)}
		+V.
		\]
	\end{definition}
	
	\begin{lemma}
		Der Operator \(H\) ist wohldefiniert auf \(\HE\).
	\end{lemma}
	
	\begin{proof}
		Die Invarianz von \(\XE\) unter \(T\) und \(F_{ij}^{(g)}\) wurde bereits gezeigt; \(V\) ist diagonal.
	\end{proof}
	
	\section{Orbitstruktur}
	
	\begin{definition}
		Für \(X\in \XE\) sei
		\[
		\mathcal O(X):=\{T^kX: k=0,1,2,3\}
		\]
		der Shift-Orbit von \(X\).
	\end{definition}
	
	\begin{lemma}
		Jeder Shift-Orbit hat Länge \(1\), \(2\) oder \(4\).
	\end{lemma}
	
	\begin{proof}
		Aus \(T^4=I\) folgt, dass die Orbitlänge ein Teiler von \(4\) ist.
	\end{proof}
	
	\begin{lemma}
		Für
		\[
		X_0=(A,B,C,E)
		\]
		hat der Shift-Orbit Länge \(4\):
		\[
		(A,B,C,E),\ (B,C,E,A),\ (C,E,A,B),\ (E,A,B,C).
		\]
	\end{lemma}
	
	\begin{lemma}
		Die Eigenwerte der Restriktion von \(T\) auf einen Orbit der Länge \(4\) sind
		\[
		1,\ i,\ -1,\ -i.
		\]
	\end{lemma}
	
	\begin{proof}
		Auf einem Viererorbit wirkt \(T\) als zyklische Permutationsmatrix der Länge \(4\). Deren Eigenwerte sind die vierten Einheitswurzeln.
	\end{proof}
	
	\begin{corollary}
		Für den Transportterm
		\[
		H_{\mathrm{trans}}:=-t(T+T^{-1})
		\]
		sind auf jedem Orbit der Länge \(4\) die Eigenwerte
		\[
		-2t,\ 0,\ 0,\ 2t.
		\]
	\end{corollary}
	
	\begin{proof}
		Ist \(\omega\in\{1,i,-1,-i\}\) ein Eigenwert von \(T\), so ist
		\[
		-t(\omega+\omega^{-1})
		\]
		der zugehörige Eigenwert von \(H_{\mathrm{trans}}\). Einsetzen liefert die Behauptung.
	\end{proof}
	
	\section{Typklassen neutraler Vierlinge}
	
	\begin{definition}
		Ein Zustand \(X\in \XE\) heißt vom Typ \emph{EEEE}, falls
		\[
		X=(E,E,E,E).
		\]
	\end{definition}
	
	\begin{definition}
		Ein Zustand \(X\in \XE\) heißt vom Typ \emph{AABB}, falls seine Multimenge aus vier Einträgen zwei verschiedene Elemente mit Vielfachheit \(2\) enthält.
	\end{definition}
	
	\begin{definition}
		Ein Zustand \(X\in \XE\) heißt vom Typ \emph{ABCE}, falls jedes der vier Elemente \(E,A,B,C\) genau einmal vorkommt.
	\end{definition}
	
	\begin{remark}
		Die Typen AABB und ABCE sind disjunkt. Beide liegen vollständig im neutralen Sektor \(\XE\).
	\end{remark}
	
	\section{Numerische Parametrisierung}
	
	Für die numerischen Rechnungen wird das Potential in der Form
	\[
	v(X)=\mu\,N_{\mathrm{nt}}(X)+\alpha\,I_{ABCE}(X)
	\]
	gewählt, wobei
	\begin{itemize}
		\item \(N_{\mathrm{nt}}(X)\) die Anzahl der nichttrivialen Einträge in \(X\) bezeichnet,
		\item \(I_{ABCE}(X)\in\{0,1\}\) die Indikatorfunktion des Typs ABCE ist.
	\end{itemize}
	
	Im Folgenden werden
	\[
	t=1,\qquad \lambda=0.15,\qquad \mu=0.4
	\]
	fixiert, während \(\alpha\) variiert wird.
	
	\section{Numerische Resultate}
	
	\begin{definition}
		Für einen normierten Eigenzustand \(\psi=\sum_{X\in \XE} c_X |X\rangle\) und eine Typklasse \(\mathcal T\subseteq \XE\) definieren wir ihr Gewicht durch
		\[
		\wt_{\mathcal T}(\psi):=\sum_{X\in\mathcal T}|c_X|^2.
		\]
	\end{definition}
	
	\subsection{Grundzustand bei \texorpdfstring{$\alpha=-0.3$}{alpha=-0.3}}
	
	Für \(\alpha=-0.3\) erhält man numerisch
	\[
	E_0\approx -3.68503172,\qquad E_1\approx -2.22057451.
	\]
	Die Gewichte des Grundzustands \(\psi_0\) in den drei Typklassen sind
	\[
	\wt_{AABB}(\psi_0)\approx 0.37584705,\qquad
	\wt_{ABCE}(\psi_0)\approx 0.17788503,\qquad
	\wt_{EEEE}(\psi_0)\approx 0.05361052.
	\]
	Für den ersten angeregten Zustand \(\psi_1\) gilt
	\[
	\wt_{EEEE}(\psi_1)\approx 0.67243007.
	\]
	
	\begin{proposition}
		Für \(\alpha=-0.3\) ist der Grundzustand \(\psi_0\) nicht vom Typ EEEE dominiert. Vielmehr gilt
		\[
		\wt_{AABB}(\psi_0)>\wt_{ABCE}(\psi_0)>\wt_{EEEE}(\psi_0).
		\]
	\end{proposition}
	
	\begin{proof}
		Dies ist unmittelbare Konsequenz der numerisch bestimmten Gewichte.
	\end{proof}
	
	\subsection{Scan in \texorpdfstring{$\alpha$}{alpha}}
	
	\begin{table}[h]
		\centering
		\begin{tabular}{rcccc}
			\toprule
			$\alpha$ & $E_0$ & AABB & ABCE & EEEE \\
			\midrule
			$-1.00$ & $-4.00471086$ & $0.427823$ & $0.526006$ & $0.029035$ \\
			$-0.75$ & $-3.87939297$ & $0.467105$ & $0.476415$ & $0.036655$ \\
			$-0.50$ & $-3.76650374$ & $0.505010$ & $0.426820$ & $0.045587$ \\
			$-0.25$ & $-3.66586282$ & $0.540303$ & $0.378661$ & $0.055726$ \\
			$0.00$  & $-3.57694646$ & $0.572016$ & $0.333212$ & $0.066867$ \\
			$0.25$  & $-3.49895031$ & $0.599557$ & $0.291426$ & $0.078731$ \\
			$0.50$  & $-3.43088069$ & $0.622728$ & $0.253861$ & $0.091008$ \\
			\bottomrule
		\end{tabular}
		\caption{Gewichte des Grundzustands in den Typklassen AABB, ABCE und EEEE.}
	\end{table}
	
	Ein Feinscan im relevanten Bereich ergibt:
	
	\begin{table}[h]
		\centering
		\begin{tabular}{rccccc}
			\toprule
			$\alpha$ & $E_0$ & AABB & ABCE & EEEE & ABCE$-$AABB \\
			\midrule
			$-0.80$ & $-3.90346289$ & $0.459317$ & $0.486381$ & $0.035024$ & $+0.027064$ \\
			$-0.70$ & $-3.85582136$ & $0.474838$ & $0.466450$ & $0.038338$ & $-0.008389$ \\
			$-0.60$ & $-3.81017142$ & $0.490097$ & $0.446565$ & $0.041861$ & $-0.043532$ \\
			$-0.50$ & $-3.76650374$ & $0.505010$ & $0.426820$ & $0.045587$ & $-0.078190$ \\
			$-0.40$ & $-3.72479970$ & $0.519502$ & $0.407307$ & $0.049507$ & $-0.112195$ \\
			$-0.30$ & $-3.68503172$ & $0.533505$ & $0.388113$ & $0.053611$ & $-0.145392$ \\
			$-0.20$ & $-3.64716380$ & $0.546957$ & $0.369319$ & $0.057883$ & $-0.177638$ \\
			$-0.10$ & $-3.61115223$ & $0.559808$ & $0.350997$ & $0.062307$ & $-0.208811$ \\
			$0.00$  & $-3.57694646$ & $0.572016$ & $0.333212$ & $0.066867$ & $-0.238804$ \\
			$0.10$  & $-3.54449005$ & $0.583553$ & $0.316018$ & $0.071543$ & $-0.267535$ \\
			\bottomrule
		\end{tabular}
		\caption{Feinscan des Umschaltbereichs.}
	\end{table}
	
	\begin{proposition}
		Es existiert ein numerischer Umschaltbereich, in dem die Dominanz des Grundzustands von der Typklasse ABCE auf die Typklasse AABB übergeht. Für die gegebene Parameterwahl liegt der Umschlag bei ungefähr
		\[
		\alpha_c\approx -0.72.
		\]
	\end{proposition}
	
	\begin{proof}
		Aus der Tabelle folgt
		\[
		\wt_{ABCE}(\psi_0)-\wt_{AABB}(\psi_0)>0
		\quad\text{für }\alpha=-0.80,
		\]
		aber
		\[
		\wt_{ABCE}(\psi_0)-\wt_{AABB}(\psi_0)<0
		\quad\text{für }\alpha=-0.70.
		\]
		Damit liegt der Vorzeichenwechsel zwischen diesen beiden Werten. Lineare Interpolation liefert \(\alpha_c\approx -0.724\).
	\end{proof}
	
	\begin{corollary}
		Im untersuchten Parameterbereich bleibt die Typklasse EEEE im Grundzustand untergeordnet.
	\end{corollary}
	
	\begin{proof}
		Dies folgt unmittelbar aus den in den Tabellen angegebenen Gewichten.
	\end{proof}
	
	\section{Schlussbemerkung}
	
	Der diskrete Vierlingsoperator auf \(\HE\) besitzt bereits in endlicher Dimension eine nichttriviale Grundzustandsstruktur. Die numerische Analyse zeigt insbesondere:
	\begin{enumerate}[label=(\roman*)]
		\item der Grundzustand ist nicht trivial,
		\item die Typklassen AABB und ABCE konkurrieren im Grundzustand,
		\item der Steuerparameter \(\alpha\) erzeugt einen numerisch lokalisierbaren Umschlag,
		\item die Typklasse EEEE bleibt im untersuchten Bereich sekundär.
	\end{enumerate}
	
	Damit liegt ein explizites endliches Modell vor, in dem verschiedene Formen neutraler Vierlingsordnung spektral unterschieden werden.
	
	\section*{Ausblick}
	
	Naheliegende Folgeschritte sind:
	\begin{enumerate}[label=(\arabic*)]
		\item ein zweidimensionaler Scan in \((\alpha,\lambda)\),
		\item die vollständige Zerlegung von \(\HE\) in Shift-Orbits,
		\item eine systematische Untersuchung der Entartungen,
		\item die Einführung weiterer diagonaler oder orbitabhängiger Potentiale.
	\end{enumerate}
	
	\begin{thebibliography}{9}
		
		\bibitem{ConwaySmith}
		J. H. Conway, D. A. Smith,
		\textit{On Quaternions and Octonions},
		A K Peters, 2003.
		
		\bibitem{Hall}
		K. v. Klitzing, G. Dorda, M. Pepper,
		\textit{New Method for High-Accuracy Determination of the Fine-Structure Constant Based on Quantized Hall Resistance},
		Phys. Rev. Lett. \textbf{45} (1980), 494--497.
		
		\bibitem{AB}
		Y. Aharonov, D. Bohm,
		\textit{Significance of Electromagnetic Potentials in the Quantum Theory},
		Phys. Rev. \textbf{115} (1959), 485--491.
		
	\end{thebibliography}
	
\end{document}